\section{A cubic lower-bound family}\label{sec:cubic-family}

\begin{theorem}\label{thm:cubic-family}
There is an absolute constant $C$ such that, for every real $X\ge2$,
\[
F(X)\ge \frac{4^{1/3}}{3}X^{1/3}-C.
\]
\end{theorem}

\begin{proof}
Erd\H{o}s posed the product-support counting problem in
\cite{Erdos1980}. We know of no published source for this particular
polynomial family.
Take integers $t\ge7$ satisfying
\[
t\not\equiv0\pmod4,\qquad t\equiv1,2,7,8\pmod9,
\]
and put
\[
m=\frac{t^2-4}{3},\qquad n=\frac{t^3-3t-2}{4}.
\]
The resulting integers are positive, and
\[
12(n-m)=3t^3-4t^2-9t+10>0
\]
for $t\ge7$, so $m<n$. Direct multiplication gives
\[
m(m+1)=\frac{(t^2-4)(t^2-1)}9,\qquad
n(n+1)=\frac{(t^2-4)(t^2-1)^2}{16}.
\]
For every prime other than $2$ and $3$, the two products have the same
prime support. If $t\equiv\pm1\pmod9$, then $9\mid t^2-1$ and
$3\mid t^2-4$; if $t\equiv\pm2\pmod9$, then $9\mid t^2-4$ and
$3\mid t^2-1$. Thus division by $9$ leaves a factor $3$ in the first
product, and the second product is also divisible by $3$.

If $t$ is odd, then $8\mid t^2-1$, so both products remain even after
their displayed divisions. If $t\equiv2\pmod4$, then
$32\mid t^2-4$, which again leaves both products even after division by
$16$. Hence the two products have equal radicals.

The conditions specify twelve residue classes modulo $36$, of density
$1/3$. Since $n=t^3/4+O(t)$, the number of such $t$ with $n\le X$ is
$(4^{1/3}/3)X^{1/3}+O(1)$. The map $t\mapsto m$ is strictly increasing
for $t\ge7$, so these give distinct pairs counted by $F(X)$.
\end{proof}
